% ========================================
% Chapitre 4 : Exemples de modèles à une période
% ========================================
\chapter{Exemples de modèles à une période}

\section{Le modèle binomial}
On se place dans un marché à une période $T=1$ avec un univers $\Omega = \{\omega_u, \omega_d\}$. La probabilité historique est $P(\omega_u)=p > 0$ et $P(\omega_d)=1-p$.
\begin{itemize}
    \item Un actif sans risque dont le prix actualisé est constant : $S_0^0 = 1, S_1^0 = e^r$. Le taux d'intérêt est $r > 0$.
    \item Un actif risqué avec un prix initial $S_0^1=s$ et des prix finaux $S_1^1(\omega_u) = su$ et $S_1^1(\omega_d) = sd$, avec $u > d$.
\end{itemize}

\begin{proposition}
L'hypothèse de non-arbitrage (NA) est satisfaite si et seulement si $d < e^r < u$.
Dans ce cas, il existe une unique mesure de probabilité martingale équivalente $Q \in \mathcal{M}^e(S)$, définie par :
\begin{align*}
    q = Q(\omega_u) = \frac{e^r - d}{u - d} \\
    1-q = Q(\omega_d) = \frac{u - e^r}{u - d}
\end{align*}
Puisque $\mathcal{M}^e(S)$ est un singleton, le marché est complet.
\end{proposition}

\begin{proof}
(NA) $\iff \mathcal{M}^e(S) \neq \emptyset$. On cherche $Q$ équivalente à $P$ telle que $E_Q[\hat{S}_1^1] = \hat{S}_0^1$, où $\hat{S}$ est le prix actualisé.
\begin{equation}
E_Q\left[\frac{S_1^1}{S_1^0}\right] = \frac{S_0^1}{S_0^0} \iff q \frac{su}{e^r} + (1-q) \frac{sd}{e^r} = s
\end{equation}
En simplifiant par $s$ et en multipliant par $e^r$, on obtient :
\begin{equation}
q u + (1-q) d = e^r \iff q(u-d) = e^r - d
\end{equation}
Cette équation a une solution unique $q = \frac{e^r - d}{u - d}$. Les conditions d'équivalence $q>0$ et $1-q>0$ sont équivalentes à $e^r-d>0$ et $u-e^r>0$, soit $d < e^r < u$.
\end{proof}

\subsection{Prix et couverture d'un actif contingent}
Considérons un actif contingent $\hat{B}$ dont le payoff à l'échéance est $\hat{B}(\omega_u) = \hat{B}_u$ et $\hat{B}(\omega_d) = \hat{B}_d$.
Le marché étant complet, il existe un prix unique sans arbitrage $A_0$ et un portefeuille de réplication unique $(H_0^0, H_0^1)$.

\paragraph{Prix sans arbitrage :}
Le prix est donné par l'espérance du payoff actualisé sous la mesure martingale $Q$ :
\begin{equation}
A_0 = E_Q\left[\frac{\hat{B}}{e^r}\right] = \frac{1}{e^r} \left( q \hat{B}_u + (1-q) \hat{B}_d \right) \quad \text{avec} \quad q = \frac{e^r - d}{u - d}.
\end{equation}

\paragraph{Portefeuille de couverture :}
Le portefeuille $(H_0^0, H_0^1)$ doit satisfaire $A_0 = H_0^0 S_0^0 + H_0^1 S_0^1 = H_0^0 + H_0^1 s$.
La valeur du portefeuille à l'instant $T=1$ doit répliquer le payoff :
\begin{equation}
H_0^0 S_1^0 + H_0^1 S_1^1 = \hat{B}
\end{equation}
Ceci conduit au système d'équations :
\begin{equation}
\begin{cases}
    H_0^0 e^r + H_0^1 su = \hat{B}_u \\
    H_0^0 e^r + H_0^1 sd = \hat{B}_d
\end{cases}
\end{equation}
La soustraction des deux lignes donne $H_0^1 s(u-d) = \hat{B}_u - \hat{B}_d$. On en déduit la quantité d'actif risqué à détenir (appelée le "Delta") :
\begin{equation}
H_0^1 = \frac{\hat{B}_u - \hat{B}_d}{s(u-d)}
\end{equation}
Et la quantité d'actif sans risque :
\begin{equation}
H_0^0 = A_0 - H_0^1 s
\end{equation}
Dans le modèle binomial avec $d < e^r < u$, je peux pricer et couvrir explicitement n'importe quel actif contingent.

\section{Le modèle trinomial}
On étend le modèle précédent à un univers à trois états $\Omega = \{\omega_u, \omega_m, \omega_d\}$. Les prix finaux de l'actif risqué sont $su, sm, sd$ avec $u > m > d$.

\begin{proposition}
L'hypothèse de non-arbitrage (NA) est satisfaite si et seulement si $d < e^r < u$.
\end{proposition}

\begin{proof}

(NA) $\iff \mathcal{M}^e(S) \neq \emptyset$. On cherche une probabilité $Q=(q_1, q_2, q_3)$ avec $q_i = Q(\omega_i) > 0$ telle que :
\begin{align}
    q_1 + q_2 + q_3 &= 1 \quad (P_1) \text{ (condition de probabilité)} \\
    q_1 u + q_2 m + q_3 d &= e^r \quad (P_2) \text{ (condition de martingale)}
\end{align}

\begin{tcolorbox}[title=Visualisation géométrique de l'arbitrage, colback=yellow!10!white, colframe=yellow!50!black]
L'ensemble des probabilités est représenté par un \textbf{simplexe} (triangle rouge). Pour qu'il n'y ait pas d'arbitrage (NA), le plan de martingale $(P_2)$ doit impérativement traverser l'intérieur de ce triangle. Cela se produit si et seulement si $d < e^r < u$.
\end{tcolorbox}
\end{proof}

\paragraph{Comment lire la Figure 4.1 :}
Cette figure illustre la structure de l'ensemble des mesures martingales équivalentes $\mathcal{M}^e(S)$. Plusieurs éléments clés sont à observer :
\begin{itemize}
    \item \textbf{Le Simplexe (Rouge) :} Il contient toutes les probabilités $Q$ telles que $\sum q_i = 1$. L'intérieur représente les probabilités \textit{équivalentes} ($q_i > 0$).
    \item \textbf{L'Ensemble Martingale ($\mathcal{M}^e(S)$) :} C'est le segment de droite (bleu ou vert) résultant de l'intersection du simplexe avec le plan de martingale.
    \item \textbf{Le Carré en pointillés (Gris) :} Sur le plan horizontal ($q_1, q_3$), il délimite le domaine $[0,1] \times [0,1]$. Il représente l'espace où chaque probabilité individuelle est valide. Le simplexe n'occupe que la partie "basse" ($q_1+q_3 \le 1$) car la somme des trois probabilités doit être exactement 1.
    \item \textbf{Dépendance aux paramètres :} Selon que le rendement "du milieu" $m$ est supérieur ou inférieur au taux sans risque $e^r$, le segment pivote pour relier différents axes du simplexe.
\end{itemize}

\begin{figure}[H]
\centering
\tdplotsetmaincoords{70}{120}
\begin{tikzpicture}[tdplot_main_coords, scale=4]
    \coordinate (O) at (0,0,0);
    
    % Axes setup:
    % q2 Up -> Z axis (0,0,1)
    % q1 Right -> Y axis (0,1,0)
    % q3 Down/Left -> X axis (1,0,0)
    
    \draw[-{Latex[length=2mm]}] (O) -- (1.5,0,0) node[anchor=north east]{$q_3$};
    \draw[-{Latex[length=2mm]}] (O) -- (0,1.5,0) node[anchor=north west]{$q_1$};
    \draw[-{Latex[length=2mm]}] (O) -- (0,0,1.5) node[anchor=south]{$q_2$};
    
    % Unit Square on floor (q1, q3)
    \draw[dashed, gray!80] (1,0,0) -- (1,1,0) -- (0,1,0);
    \node[gray, font=\tiny, anchor=south west] at (1,1,0) {$(1,1)$};

    % Vertices of Simplex on axes (adjusted to match professor's schema)
    \coordinate (Q3) at (1.25,0,0); % q3 axis intercept - outside unit square
    \coordinate (Q1) at (0,0.75,0); % q1 axis intercept - inside unit square
    \coordinate (Q2) at (0,0,1); % q2 axis intercept

    % Simplex (P1) - RED
    \draw[fill=red!5, opacity=0.6] (Q1) -- (Q2) -- (Q3) -- cycle;
    \draw[red, thick] (Q1) -- (Q2);
    \draw[red, thick] (Q2) -- (Q3);
    \draw[red, thick] (Q3) -- (Q1);
    
    \node[red, font=\bfseries] at (0.2, 0.2, 0.9) {Simplex $(P_1)$};

    % Blue Setup: Line on Q1-Q3 edge to Q2-Q3 edge
    % Point on q1-q3 edge (floor): parametric from (0,0.75,0) to (1.25,0,0)
    \coordinate (Blue_P1) at (0.9, 0.2, 0); 
    % Point on q2-q3 edge (x-z plane): parametric from (0,0,1) to (1.25,0,0)
    \coordinate (Blue_P2) at (0.85, 0, 0.3);
    
    \draw[ultra thick, blue] (Blue_P1) -- (Blue_P2) node[pos=0.1, left, xshift=-10pt, yshift=5pt, font=\scriptsize]{$\mathcal{M}^e(S)$};
    \fill[blue] (Blue_P1) circle (0.4pt);
    \fill[blue] (Blue_P2) circle (0.4pt);
    
    % Blue Normal n1
    \draw[->, blue, thick] (O) -- (1.5, 1.5, 1.5) node[above right, blue, inner sep=2pt] {$\vec{n}_1$};

    % Green Setup: Line on Q1-Q3 edge to Q1-Q2 edge
    % Point on q1-q3 edge (floor): closer to Q1
    \coordinate (Green_P1) at (0.4, 0.5, 0);
    % Point on q1-q2 edge (y-z plane): parametric from (0,0.75,0) to (0,0,1)
    \coordinate (Green_P2) at (0, 0.4, 0.45);

    \draw[ultra thick, green!60!black] (Green_P1) -- (Green_P2) node[pos=0.1, right, xshift=10pt, yshift=5pt, font=\scriptsize]{$\mathcal{M}^e(S)$};
    \fill[green!60!black] (Green_P1) circle (0.4pt);
    \fill[green!60!black] (Green_P2) circle (0.4pt);
    
    % Green Normal n2
    \draw[->, green!60!black, thick] (O) -- (0.2, 0.9, 0.4) node[right, green!60!black, inner sep=3pt] {$\vec{n}_2$};

    % Labels for points
    \node[below=8pt, blue] at (Blue_P1) {$\bar{P}_1$}; 
    \node[left=10pt, blue] at (Blue_P2) {$\bar{P}_2$};
    
    \node[below right=8pt, green!60!black] at (Green_P1) {$\bar{P}_1$}; 
    \node[right=10pt, green!60!black] at (Green_P2) {$\bar{P}_2'$};

\end{tikzpicture}
\caption{Géométrie de l'incomplétude dans le modèle trinomial. L'intersection (segment épais) entre l'espace des probabilités et la contrainte de martingale montre l'existence d'une infinité de mesures RNM. La configuration bleue correspond à $m \ge e^r$ et la verte à $m < e^r$.}
\end{figure}

\begin{tcolorbox}[title=L'intérêt de la vision géométrique, colback=red!5!white, colframe=red!50!black]
Cette vision permet de comprendre que le prix d'un actif n'est plus unique : il peut varier le long du segment $\mathcal{M}^e(S)$. Les bornes de prix sans arbitrage sont simplement les valeurs de l'actif aux extrémités du segment ($\bar{P}_1$ et $\bar{P}_2$).
\end{tcolorbox}

Plus généralement, pour obtenir $\mathcal{M}^e(S)$, on cherche l'intersection entre l'intérieur du simplexe (ouvert) et le plan $(P_2)$ d'équation :
\begin{equation}
    q_1 u + q_2 m + q_3 d = e^r
\end{equation}
Le vecteur suivant est orthogonal au plan $(P_2)$ :
\begin{equation}
    \vec{V} = \begin{pmatrix} u/e^r \\ m/e^r \\ d/e^r \end{pmatrix}
\end{equation}
Sous l'hypothèse de non-arbitrage (NA), on a toujours :
\begin{equation}
    \frac{u}{e^r} > 1 \quad \text{et} \quad 0 < \frac{d}{e^r} < 1
\end{equation}
L'intersection dépend alors de la position de $\frac{m}{e^r}$ par rapport à 1. On distingue deux cas (illustrés sur la Figure 4.1) :
Pour trouver les bornes de prix d'une créance, il suffit d'évaluer son espérance aux deux points extrêmes de ce segment, notés $\bar{P}_1$ et $\bar{P}_2$.

La forme exacte de ces points extrêmes dépend de la position du rendement intermédiaire $m$ par rapport au taux sans risque $e^r$.
L'ensemble des mesures martingales $\mathcal{M}^e(S)$ est l'intersection de deux contraintes géométriques : le simplexe des probabilités ($q_1+q_2+q_3=1, q_i>0$) et le plan de la condition de martingale ($q_1u + q_2m + q_3d = e^r$). Dans le cas trinomial avec $d < e^r < u$, cette intersection est un segment de droite. Les bornes des prix sans arbitrage sont atteintes aux extrémités de ce segment. Ces points extrêmes correspondent aux mesures martingales qui "ignorent" l'un des trois états possibles, c'est-à-dire en lui attribuant une probabilité nulle ($q_i=0$).

Détaillons le calcul pour le cas où $m \ge e^r$.

\begin{itemize}
    \item \textbf{Calcul de $\bar{P}_1$ :} Ce point est obtenu en supposant que l'état intermédiaire $m$ a une probabilité nulle, soit $q_2=0$. Le système d'équations se réduit à :
    \begin{enumerate}
        \item $q_1 + q_3 = 1$
        \item $q_1u + q_3d = e^r$
    \end{enumerate}
    En substituant $q_3 = 1 - q_1$ dans la seconde équation, on obtient $q_1u + (1 - q_1)d = e^r$, ce qui donne $q_1(u - d) = e^r - d$. La solution est :
    \[
    q_1 = \frac{e^r - d}{u - d} \quad \text{et} \quad q_3 = 1 - q_1 = \frac{u - e^r}{u - d}.
    \]
    On a donc $\bar{P}_1 = ( \frac{e^r-d}{u-d}, 0, \frac{u-e^r}{u-d} )$.

    \item \textbf{Calcul de $\bar{P}_2$ :} Ce point est obtenu en posant $q_1=0$. Le système devient :
    \begin{enumerate}
        \item $q_2 + q_3 = 1$
        \item $q_2m + q_3d = e^r$
    \end{enumerate}
    La résolution est analogue et donne :
    \[
    q_2 = \frac{e^r - d}{m - d} \quad \text{et} \quad q_3 = \frac{m - e^r}{m - d}.
    \]
    On a donc $\bar{P}_2 = ( 0, \frac{e^r-d}{m-d}, \frac{m-e^r}{m-d} )$.
\end{itemize}

Le calcul pour le cas $m < e^r$ suit la même logique, en posant $q_2=0$ puis $q_3=0$.

\subsection{Prix des actifs contingents}
Pour un actif contingent $\hat{B}$ avec des payoffs $(\hat{B}_u, \hat{B}_m, \hat{B}_d)$, l'ensemble des prix sans arbitrage est l'intervalle ouvert $]\underline{\Pi}(\hat{B}), \overline{\Pi}(\hat{B})[$, où :
\begin{itemize}
    \item $\underline{\Pi}(\hat{B}) = \inf_{Q \in \mathcal{M}^e(S)} E_Q[\hat{B}/e^r]$
    \item $\overline{\Pi}(\hat{B}) = \sup_{Q \in \mathcal{M}^e(S)} E_Q[\hat{B}/e^r]$
\end{itemize}
L'application $Q \mapsto E_Q[\cdot]$ étant linéaire et l'ensemble $\mathcal{M}^e(S)$ étant convexe (un segment), le supremum et l'infimum sont atteints aux points extrêmes du segment. Ces points extrêmes, notés $\bar{P}_1$ et $\bar{P}_2$, correspondent aux cas où l'un des états a une probabilité nulle.

Deux cas se présentent en fonction de la position de $m$ par rapport à $e^r$ :
\begin{enumerate}
    \item Si $m \ge e^r$ : Les points extrêmes sont $\bar{P}_1$ (intersection avec $q_2=0$) et $\bar{P}_2$ (intersection avec $q_1=0$).
        \begin{itemize}
            \item $\bar{P}_1 = (\frac{e^r-d}{u-d}, 0, \frac{u-e^r}{u-d})$
            \item $\bar{P}_2 = (0, \frac{e^r-d}{m-d}, \frac{m-e^r}{m-d})$
        \end{itemize}
    \item Si $m < e^r$ : Les points extrêmes sont $\bar{P}_1$ (intersection avec $q_2=0$) et $\bar{P}_2$ (intersection avec $q_3=0$).
        \begin{itemize}
            \item $\bar{P}_1 = (\frac{e^r-d}{u-d}, 0, \frac{u-e^r}{u-d})$
            \item $\bar{P}_2 = (\frac{e^r-m}{u-m}, \frac{u-e^r}{u-m}, 0)$
        \end{itemize}
\end{enumerate}

\begin{exemple}[Option d'achat européenne]
Considérons une option d'achat (call) de strike $K$ sur l'actif risqué, de sorte que le payoff actualisé soit $\hat{B} = (S_1^1 - K)^+/e^r$. Supposons $su > K > sm$ et $sd < K$. Les payoffs sont :
\[
\hat{B}(\omega_u) = \frac{su-K}{e^r}, \quad \hat{B}(\omega_m) = 0, \quad \hat{B}(\omega_d) = 0
\]
Les bornes de prix sont alors :
\begin{itemize}
    \item $\overline{\Pi}(\hat{B}) = \max(E_{\bar{P}_1}[\hat{B}/e^r], E_{\bar{P}_2}[\hat{B}/e^r])$.
    Comme $u$ est le seul état où le payoff est non-nul, le sup est atteint pour la probabilité qui maximise $q_1$. C'est $\bar{P}_1$ si $m < e^r$ ou $\bar{P}_2$ si $m > e^r$ (selon les valeurs).
    Supposons $m < e^r$. Le max de $q_1$ est en $\bar{P}_2$.
    $\overline{\Pi}(\hat{B}) = E_{\bar{P}_2}[\hat{B}/e^r] = \frac{e^r-m}{u-m} \cdot \frac{su-K}{e^r}$.
    \item $\underline{\Pi}(\hat{B}) = \min(E_{\bar{P}_1}[\hat{B}/e^r], E_{\bar{P}_2}[\hat{B}/e^r])$.
    Le min est en $\bar{P}_1$.
    $\underline{\Pi}(\hat{B}) = E_{\bar{P}_1}[\hat{B}/e^r] = \frac{e^r-d}{u-d} \cdot \frac{su-K}{e^r}$.
\end{itemize}
\end{exemple}

\begin{exemple}[Exercice numérique - Modèle trinomial]
Considérons un modèle trinomial à une période avec les paramètres suivants : taux sans risque $r=0$ (donc $e^r = 1$), prix initial $S_0^1 = 1$, facteurs de rendement $u=1.2$, $m=1$, $d=0.8$. On s'intéresse à une option d'achat de strike $K = 0.9$, dont le payoff actualisé est $\hat{B} = (\hat{S}^1_1 - 0.9)^+$.

\textbf{Questions :}
\begin{enumerate}
    \item Vérifier la condition de non-arbitrage.
    \item Calculer les payoffs dans chaque état.
    \item Déterminer les mesures martingales extrêmes.
    \item En déduire l'intervalle des prix sans arbitrage.
    \item Construire un portefeuille de sur-réplication.
\end{enumerate}
\end{exemple}

\section*{Solution de l'Exercice Numérique 4.3}

Cette sous-section propose une solution détaillée de l'exercice ci-dessus.

\begin{enumerate}
    \item \textbf{Paramètres :}
    \begin{itemize}
        \item Taux sans risque : $r=0$, donc $e^r = 1$.
        \item Prix initial : $S_0^1 = 1$.
        \item Facteurs de rendement : $u=1.2, m=1, d=0.8$.
        \item Créance contingente : $\hat{B} = (\hat{S}^1_1 - 0.9)^+$.
    \end{itemize}
    \item \textbf{Analyse du marché :} On vérifie la condition de non-arbitrage : $d < e^r < u$ devient $0.8 < 1 < 1.2$. La condition est satisfaite. Le marché est incomplet.
    \item \textbf{Calcul des payoffs :}
    \begin{itemize}
        \item État $u$ : $\hat{S}^1_1 = 1.2$, $\hat{B}_u = (1.2 - 0.9)^+ = 0.3$.
        \item État $m$ : $\hat{S}^1_1 = 1.0$, $\hat{B}_m = (1.0 - 0.9)^+ = 0.1$.
        \item État $d$ : $\hat{S}^1_1 = 0.8$, $\hat{B}_d = (0.8 - 0.9)^+ = 0$.
    \end{itemize}
    \item \textbf{Détermination des mesures extrêmes :} Nous sommes dans le cas $m = e^r = 1$, qui est un cas limite de $m \ge e^r$.
    \begin{itemize}
        \item $\bar{P}_1$ (avec $q_2=0$, donc $q_1+q_3=1$ et $1.2q_1 + 0.8q_3 = 1$) :
        \[
        q_1 = \frac{1-0.8}{1.2-0.8} = \frac{0.2}{0.4} = 0.5.
        \]
        Donc $\bar{P}_1 = (0.5, 0, 0.5)$.
        \item $\bar{P}_2$ (avec $q_1=0$, donc $q_2+q_3=1$ et $1.0q_2 + 0.8q_3 = 1$) :
        \[
        q_2 = \frac{1-0.8}{1-0.8} = 1.
        \]
        Donc $\bar{P}_2 = (0, 1, 0)$.
    \end{itemize}
    \item \textbf{Calcul des bornes de prix :} L'intervalle des prix sans arbitrage est $]\underline{\Pi}(\hat{B}), \overline{\Pi}(\hat{B})[$.
    \begin{itemize}
        \item $E_{\bar{P}_1}[\hat{B}] = 0.5 \times 0.3 + 0 \times 0.1 + 0.5 \times 0 = 0.15$.
        \item $E_{\bar{P}_2}[\hat{B}] = 0 \times 0.3 + 1 \times 0.1 + 0 \times 0 = 0.1$.
    \end{itemize}
    L'intervalle des prix est donc $]0.1, 0.15[$.
    \item \textbf{Calcul du portefeuille de sur-réplication :} Le coût minimal pour sur-répliquer est $\alpha = \overline{\Pi}(\hat{B}) = 0.15$. Il faut trouver un portefeuille $(H_0^0, H_0^1)$ tel que son coût initial soit 0.15 et sa valeur finale domine le payoff $\hat{B}$ dans tous les états.
    \begin{itemize}
        \item Coût initial : $\hat{V}_0 = H_0^0 \cdot 1 + H_0^1 \cdot 1 = 0.15$.
        \item Le portefeuille proposé est $H_0^1 = 0.75$ et $H_0^0 = -0.6$.
        \item Vérification du coût : $-0.6 + 0.75 = 0.15$. Correct.
        \item Vérification de la sur-réplication à $T=1$ :
        \begin{itemize}
            \item État $u$: $\hat{V}_1 = -0.6 \times 1 + 0.75 \times 1.2 = -0.6 + 0.9 = 0.3$. Or $\hat{B}_u = 0.3$. Donc $\hat{V}_1 \ge \hat{B}_u$. OK.
            \item État $m$: $\hat{V}_1 = -0.6 \times 1 + 0.75 \times 1.0 = 0.15$. Or $\hat{B}_m = 0.1$. Donc $\hat{V}_1 \ge \hat{B}_m$ ($0.15 \ge 0.1$). OK.
            \item État $d$: $\hat{V}_1 = -0.6 \times 1 + 0.75 \times 0.8 = -0.6 + 0.6 = 0$. Or $\hat{B}_d = 0$. Donc $\hat{V}_1 \ge \hat{B}_d$. OK.
        \end{itemize}
    \end{itemize}
    Le portefeuille sur-réplique bien le payoff pour un coût de 0.15.
\end{enumerate}
